% !TEX root = ../Thesis.tex
\section{Static Analysis}

This section is focused on tools which can be used to perform static analysis on Rust code.

\subsection{Prusti}

Prusti is a verification framework whose development has started in 2017 at the ETH Zurich \cite{astrauskas_leveraging_2019,deshmukh_prusti_2022}. It leverages Rust's type system while providing a specification language, which can be used to describe certain properties of the program for verifying \cite{astrauskas_leveraging_2019}. This verifying process is automated by translating the Rust program and specifications into the Viper \cite{jobstmann_viper_2016} intermediate language, which is used for reasoning \cite{astrauskas_leveraging_2019}. There were three main goals while designing Prusti:

\begin{itemize}
    \item Provide functional correctness properties for Rust programs, without depending on Rust's guarantees \cite{deshmukh_prusti_2022}.
    \item Make verification accessible to programmers without needing the domain knowledge of verification \cite{deshmukh_prusti_2022}.
    \item Make the usage of Prusti easy for Rust programmers \cite{deshmukh_prusti_2022}.
\end{itemize}

\subsubsection{Usability}

Usability is prioritized goal for Prusti.
Prusti integrates itself into the \textit{rustc} compiler and uses the generated \acs{mir} for further interaction \cite{deshmukh_prusti_2022}. This has the advantage that the verifier itself always stays on pair with the progress of the Rust syntax, without relying on adaptions of the tool \cite{deshmukh_prusti_2022}. Another advantage is the resulting ease of use with \textit{cargo}, which allows Prusti a convenient integration in the already present toolchain for Rust development. Furthermore, there is a extension called Prusti Assistant \cite{chair_of_programming_methodology_-_eth_zurich_prusti_nodate} for the popular code editor \textit{Visual Studio Code}. It takes care of the local setup process of Prusti and integrates itself into the editor for inline error highlighting of error messages produced by Prusti.

Prusti is fully autonomous. The MIR, generated by the compiler, gets translated \cite{astrauskas_leveraging_2019} into the Viper intermediate language, which is a part of the Viper infrastructure \cite{jobstmann_viper_2016}. The Viper infrastructure then produces the verification results and reports those back to Prusti, which then consequently reports them back to the user \cite{astrauskas_leveraging_2019}.

Prusti proposes incremental verification. Prusti makes it easy to retrofit a existing codebase with Prusti annotations. Prusti proves by default that a program will not terminate after reaching an unrecoverable state, e.g. a panic!() or unreachable!() macro call, without the need of specifications in the Rust code \cite{deshmukh_prusti_2022}. Furthermore, it can also prove the absence of integer overflows \cite{deshmukh_prusti_2022}. If desired, a developer then can add specifications to the code to get more guarantees \cite{deshmukh_prusti_2022}.

Prusti uses Rust syntax for designing specifications. However, the syntax has been extended by some keywords to help design those specifications. For example, for designing post-conditions, it uses the keyword \textit{result} to refer to the result of a given function \cite{deshmukh_prusti_2022}. Another keyword is \textit{old}, which refers to the state of a given variable, before the function has been executed \cite{deshmukh_prusti_2022}. As can be seen in \autoref{lst:prusti_pre_postcondition}, the Prusti specifications are Rust specific attributes \cite{rust_foundation_attributes_nodate}, distinguished with a Prusti keyword.

\begin{lstlisting}[language=Rust, caption=Example Rust code of a function with a pre- and post-condition for Prusti \cite{deshmukh_prusti_2022}, label={lst:prusti_pre_postcondition}]
#[requires(a.len() < usize::MAX / 2)]
#[ensures(if let Some(idx) = result { idx < a.len() && a[idx] == key } else { true })]
fn search(a: &[i32], key: i32) -> Option<usize> { /* ... */ }
\end{lstlisting}

\subsubsection{Technology}
%Näher auf Viper eingehen und Z3 und Boogie, Seperation Logic
\subsubsection{Feedback}
Due to Integration of Rustc nice error messages 
\subsubsection{Language support}
Recursion wird supported, Library calls können annotiert werden, Pledges, Contracts, Loop/Type invariants, old expressions
\subsubsection{Development}
\subsubsection{Activity}
\subsubsection{Scientific activity}
\subsubsection{Documentation}

